\section{Majority Rule}

In a paper from 1976, ``On Frustration of the Majority by Fulfilment of the Majority's Will'', the philosopher \textbf{Elizabeth Anscombe} demonstrate this potential fundamental flaw in democratic systems:

\begin{quotex}
The majority votes in the minority in a majority of cases.
\end{quotex}


\paragraph{Bloc Voting}


She created a scenario where this is clearly true; however, she assumes that each person votes as an individual independently of the other voters. That is, there is no bloc voting. However, the 2012 presidential elections in the USA shows that her thesis is even stronger in conditions of bloc voting: if smaller blocs vote overwhelmingly for candidate A, that will override the will of the majority bloc for candidate B.

Oddly enough, this is celebrated in the USA which formerly had prided itself on what was called the ``melting pot''. In a melting pot, bloc differences dissolve and roughly the same percentage of voters will hold certain positions across races, ethnicities, sexes, and so on. The opposite of the melting pot is celebrated under the name ``diversity''. Now, we don't consider diversity a value or a disvalue, since a common mind or spirit is more important to a healthy society than its composition.

Taken to the extreme, a commentator on NPR\footnote{\url{http://www.npr.org/}} last week praised the new Congress, at least the democrats, for ``looking'' more like America in its distribution. That is a view that ``intelligent'' people are expected to hold. However, just because the congress looks like Walmart on a Saturday afternoon, with the same level of cohesion, does not merit praise. If only the commentator had said something like this: ``The new congress is remarkable for the level of intelligence of the representatives, their success in practical fields, their skill in statecraft, and their renown for moral rectitude.''

Anscombe recognizes the problem with bloc voting from this real life example she provided, demonstrating that the best and the brightest are not beyond power politics:

\begin{quotex}
I know a university where representation in the authoritative assembly was made proportional to the numbers in the different departments. Large departments proceeded to forbid expansion to or even to reduce small ones, increasing themselves as every opportunity offered.
\end{quotex}


\paragraph{Stupid Decisions}


Anscombe addresses the obvious, yet seldom discussed, question of the very rationale for majority rule. It manifestly can lead to stupid decisions defined as:

\begin{quotex}
A decision is stupid if its implementation is undesirable from the point of view of the very people who wanted it. They would have it: but when they get it, it is evident that they would rather not have it; if they had only realized what it would be like for their decision to be implemented, it would not have been made.
\end{quotex}


She offers no specifics, but we could assume the implementation is too expensive, unfeasible, or has unintended consequences. She offers this alternative, even though it is usually impossible to convince the majority, except when it is too late:

\begin{quotex}
There was another decision available perhaps: if the power of making the decision had not been taken away from some particular person and vested in the majority. And it is now recognized that \emph{that} decision would have been better.
\end{quotex}


The clear, and unfortunate conclusion, is this:

\begin{quotex}
The superiority of the majority decision cannot reside in its necessarily superior wisdom. A superiority, even when the decision is stupid, is rather thought to be in this: participation in the decision-making is itself valuable; or perhaps is itself a value.
\end{quotex}


She calls this theory ``mystical'' in the sense that it is based on neither fact nor principles. She speculates that majority rule engenders a ``we'' feeling. Such a ``we'' feeling is quite powerful, shown especially in the recent elections. People had such a positive emotional reaction when ``their'' candidate won, that is overruled all rationality and even knowledge or anticipation of the likely consequences of that victory became immaterial.

\paragraph{Fairness}


Another justification for majority rule is the issue of fairness. She provides the example of people on a bus deciding where to go. I can illustrate this with some vacation goers, some of whom want to go to Busch Gardens and others to Universal Studio. It seems fair to most people that the majority's choice should prevail even if the minority is disappointed.

This assumes that the decision is collective and binding on all. An alternative may be something like a round robin whereby each person in the group takes turns deciding the activity. I can expand on an example she provides, in which, say, a family cannot agree on which TV program to watch at night. Hence, they agree to take turns.

One night, the father chooses the military channel, and everyone else has to watch it. The next night, the mother chooses a romance novel on the lifetime network. The young daughter chooses reruns of Dora the Explorer\footnote{\url{https://www.gornahoor.net/?p=1104}} on her night, again boring everyone else to death. A problem arises when the teenage son chooses to watch soft-core pornography on cable. By the standards given, his choice seems fair because it assumes there is no higher value; nevertheless, from a higher point of view, it is unfair because it has more deleterious effects on the family than simple boredom. Hence, the father needs to step in, unless, of course, he is a modern father.

\paragraph{Coercion}


Anscombe brings up the issue of coercion since the majority can vote to prohibit the minority from engaging in activity X. This decision is often arbitrary, since some men will always desire to control other men. Nevertheless, as the family TV example shows, sometimes activity X is to be banned not on the grounds of a majority decision, but rather because it may be immoral, unhealthy, or simply stupid. But this decision presupposed a common agreement on the principles to determine morality, health, and intelligence. That is not possible in a mere collection of individuals, united by no higher principle, but only by a commitment to majority rule and some form of fairness as defined above. This is the situation in the USA today.

Another point not brought up by Anscombe is the issue of finances. For example, suppose the majority on the vacation bus discovers that not only do they decide the destination, but they can extract the funds to pay for it from the minority. In that case, they may decide to travel to New York and take in a Broadway show.

Anscombe concludes with a demonstration of the possibility of tyranny even in a democracy, which proof is left to the reader. She writes:

\begin{quotex}
It will be possible for the tyrant to damage the interests of anyone or any group (that does not support him, say) while truthfully claiming democratic support for his measures.
\end{quotex}



\flrightit{Posted on 2013-01-08 by Cologero}

\begin{center}* * *\end{center}

\begin{footnotesize}
\begin{sffamily}
\texttt{Logres on 2013-01-09 at 22:45 said:}

Baron Ledhin points out that the Russians were being perfectly honest in terming their system of soviet rule a ``true democracy'' (the ``truest'', from their point of view): he cites Harold Laski (not exactly a traditionalist) who argues that parliamentary democracy can only possible work towards justice when there is an agreement on the fundamentals that provides the cohesiveness to work through various plans to implement a common goal. TV creates ``majority rule'' in America -- hence Lippman's focus on propaganda after WWII, at the starting peak of American hubris. Colonel House was instrumental in working W. Wilson up into ``crusade for democracy'' (which began at home, and still had a populist basis). And here we are today. Thanks for a very insightful article: a fantastic point was a baseline stupidity definition: ``A majority decision is self-reductively stupid when it even its proponents later later on can see the folly of it''. RJ Rushdoony has a work on the medieval/feudal basis of the Republic which is helpful (although it has flaws): \url{http://www.amazon.com/This-Independent-Republic-Rousas-Rushdoony/dp/1879998246}

\hfill

\end{sffamily}
\end{footnotesize}

